\section{Reading acquisition costs in an actual report}
\label{sec:munich-new-contracts}

The preceding example supplied the acquisition assets and the net inflows
for two identified groups. A published insurer's accounts reveal a different
level of detail. Munich Re's 2023 report shows contracts recognised for the
first time during the year, separating concluded from acquired contracts and
nononerous from onerous business. Note 27 covers segments predominantly using
general measurement or the variable fee approach; it omits two predominantly
PAA segments. These are segment disclosures, not a list of individual groups
\citep[printed pp. 252--254]{munichre2023annual}.

Consider the ERGO Life and Health Germany segment, identified in the paragraph
preceding its table. Its newly recognised contracts combine GMM and VFA
measurements. The following two columns reproduce its concluded-contract
amounts in millions of euros. Preserve the report's signs: inflows are positive,
outflows and the margin are negative. Its onerous CSM cell is blank; a dash
below preserves that blank instead of assigning a reported zero
\citep[printed pp. 252--253]{munichre2023annual}.

\begin{center}
\begin{tabular}{p{0.58\textwidth}rr}
\toprule
 & Nononerous & Onerous \\
\midrule
Present value of future inflows & 2,842 & 17 \\
Future claims, expenses and investment components & $-2,438$ & $-12$ \\
Future acquisition costs & $-243$ & $-25$ \\
Risk adjustment & $-8$ & 0 \\
Contractual service margin & $-154$ & --- \\
Reported total & 0 & $-21$ \\
\bottomrule
\end{tabular}
\end{center}

The nononerous amounts illustrate how expected future profit is deferred in
CSM. The onerous column retains a loss-bearing population separately. Munich
Re describes recognising an expected initial loss in profit or loss and
keeping onerous contracts separate from profitable ones. Combining these
columns into one net profit would erase that distinction
\citep[printed pp. 165--166]{munichre2023annual}.

Adding the displayed nononerous numbers gives $-1$, although the table reports
zero. The onerous cash-flow and risk-adjustment rows give $17-12-25-0=-20$,
against a reported total of $-21$. The report states that amounts are rounded
to millions and that totals may consequently differ. A possible set of
unreported nononerous values is $2{,}842.3-2{,}437.8-242.8-7.8-153.9=0$;
each rounds to its displayed amount. Similarly,
$16.6-12.1-25.2-0.3=-21$ is consistent with the onerous column if no CSM is
assigned. Together these hypothetical values also reproduce the table's
outflow subtotals and aggregate rows after rounding. They demonstrate
consistency under nearest-million rounding; they do not recover the issuer's
unreported measurements or establish its exact rounding convention. Retain
the original figures and blank cells \citep[printed pp. 155, 253]{munichre2023annual}.

\clearpage
The acquisition-cost row appears within \emph{future cash outflows}. Its
reported total of $243+25=268$ therefore cannot simply be entered as the
previous example's asset $A$, which arose from costs already paid before the
related groups were recognised. Nor do the table's acquired-contract columns
describe how sales acquisition costs were paid. Those columns concern a
different distinction between contract populations. Their displayed zeros
retain the same rounding qualification. The selected disclosure does not
identify separate prior acquisition assets, their impairment history or their
allocation to individual groups \citep[printed pp. 252--253]{munichre2023annual}.

A third quantity appears in the segment's service revenue. For 2023, Munich
Re reports €202 million as the portion of premium relating to recovery
of acquisition costs. Its policy explains that acquisition costs enter
insurance service expense in the same amount as the recovery included in
revenue, allocated over the coverage period. This is a period service amount;
the €268 million above is an estimate of future acquisition outflows for
newly recognised contracts. Subtracting one from the other would not produce
a remaining acquisition asset. The relevant populations, dates and movements
must first be linked \citep[printed pp. 177, 253, 264]{munichre2023annual}.

The segment's liability reconciliation makes that caution concrete: it shows
€637 million in the service-expense row labelled ``Recovery of acquisition
costs'', and €978 million of acquisition costs within cash flows. Thus the
€202 million recovery in revenue does not describe the entire expense row.
The difference of €435 million is an arithmetic residual; assigning its
cause requires the relevant populations and expense treatments. Neither that
residual nor cash less expense identifies a separate acquisition asset
\citep[printed pp. 230--231]{munichre2023annual}.

The same timing distinction applies to profit. The report gives
€819 million of CSM for services provided during 2023 in this segment.
That amount concerns service during the year, whereas €154 million is
the margin attached to the newly recognised contracts in the selected table.
It is not a release from that new-business amount alone. A reconciliation
needs the opening CSM and every intervening movement. The segment's total
insurance revenue also combines measurement approaches: €7,966 million
under GMM/VFA and €1,976 million under PAA make the reported
€9,942 million \citep[printed pp. 253, 264]{munichre2023annual}.

Even within one insurer, an acquisition-cost policy needs its applicability
conditions. Munich Re reports using the PAA option to expense acquisition
costs directly for a substantial part of its property-casualty primary
insurance business. The earlier general-measurement asset example therefore
requires the relevant treatment and group facts before it can be applied to
a real portfolio \citep[printed p. 166]{munichre2023annual}. These are readings
of the retained 2023 English translation, whose German original is binding.
They identify reported aggregates and necessary further facts; the underlying
group histories and valuations remain unresolved.

\clearpage
